% GENERATED FILE. Edit canonical lessons, then run npm run generate:books.
% canonical-source-hash: fnv1a64:8d6fa7bde7cd2432
% canonical-lessons: KA-C196-lekhaka, KA-C196-nirdesaka, KA-C196-atagara, KA-C196-preksaka, KA-C196-badagi

\chapter{Writer, Director, Player, Spectator, Carpenter}
\label{ch:ka-writer-director-player-spectator-carpenter}
\hlchaptermodality{196}

\begin{chapteropening}
\textbf{By the end of this chapter:} \emph{I can say a writer, a director, a player, a spectator and a carpenter.}

Complete the last lesson of chapter 196: I can say a writer, a director, a player, a spectator and a carpenter.
\end{chapteropening}

\section[\texorpdfstring{\kn{ಲೇಖಕ} (lēkhaka)}{ಲೇಖಕ (lēkhaka)}]{\kn{ಲೇಖಕ} (lēkhaka) — a writer, an author}

\label{lesson:KA-C196-lekhaka}



\begin{quote}
Before the new one: say the Kannada for tax, then the Kannada for a bill.
\end{quote}

\subsection*{You'll want to know: \kn{ಲೇಖಕ}}
\textbf{\kn{ಲೇಖಕ}} — \emph{lēkhaka} — \textquotedblleft{}a writer, an author\textquotedblright{}.

\begin{grammarlens}[title={What you've built}]
One more word for this chapter.
\end{grammarlens}

\subsection*{Your turn}
\begin{itemize}
\raggedright
  \item \emph{Say it:} \emph{lēkhaka}
  \item \emph{Say it:} \emph{lēkhaka}, once more
  \item \emph{Say it:} say \emph{lēkhaka}
  \item \emph{Recall:} say the Kannada for tax, then the Kannada for a bill, then say \emph{lēkhaka} again
\end{itemize}

\begin{tcolorbox}[breakable,colback=teal!4,colframe=teal!35!black,title={Before you move on}]
What does \kn{ಲೇಖಕ} mean? (\textquotedblleft{}A writer, an author\textquotedblright{}.) Say it once more.
\end{tcolorbox}
\section[\texorpdfstring{\kn{ನಿರ್ದೇಶಕ} (nirdēśaka)}{ನಿರ್ದೇಶಕ (nirdēśaka)}]{\kn{ನಿರ್ದೇಶಕ} (nirdēśaka) — a director}

\label{lesson:KA-C196-nirdesaka}



\begin{quote}
Before the new one: say the Kannada for a bill, then the Kannada for a writer.
\end{quote}

\subsection*{You'll want to know: \kn{ನಿರ್ದೇಶಕ}}
\textbf{\kn{ನಿರ್ದೇಶಕ}} — \emph{nirdēśaka} — \textquotedblleft{}a director\textquotedblright{}.

\begin{grammarlens}[title={What you've built}]
One more word for this chapter.
\end{grammarlens}

\subsection*{Your turn}
\begin{itemize}
\raggedright
  \item \emph{Say it:} \emph{nirdēśaka}
  \item \emph{Say it:} \emph{nirdēśaka}, once more
  \item \emph{Say it:} say \emph{nirdēśaka}
  \item \emph{Recall:} say the Kannada for a bill, then the Kannada for a writer, then say \emph{nirdēśaka} again
\end{itemize}

\begin{tcolorbox}[breakable,colback=teal!4,colframe=teal!35!black,title={Before you move on}]
What does \kn{ನಿರ್ದೇಶಕ} mean? (\textquotedblleft{}A director\textquotedblright{}.) Say it once more.
\end{tcolorbox}
\section[\texorpdfstring{\kn{ಆಟಗಾರ} (āṭagāra)}{ಆಟಗಾರ (āṭagāra)}]{\kn{ಆಟಗಾರ} (āṭagāra) — a player}

\label{lesson:KA-C196-atagara}



\begin{quote}
Before the new one: say the Kannada for a writer, then the Kannada for a director.
\end{quote}

\subsection*{You'll want to know: \kn{ಆಟಗಾರ}}
\textbf{\kn{ಆಟಗಾರ}} — \emph{āṭagāra} — \textquotedblleft{}a player\textquotedblright{}.

\begin{grammarlens}[title={What you've built}]
One more word for this chapter.
\end{grammarlens}

\subsection*{Your turn}
\begin{itemize}
\raggedright
  \item \emph{Say it:} \emph{āṭagāra}
  \item \emph{Say it:} \emph{āṭagāra}, once more
  \item \emph{Say it:} say \emph{āṭagāra}
  \item \emph{Recall:} say the Kannada for a writer, then the Kannada for a director, then say \emph{āṭagāra} again
\end{itemize}

\begin{tcolorbox}[breakable,colback=teal!4,colframe=teal!35!black,title={Before you move on}]
What does \kn{ಆಟಗಾರ} mean? (\textquotedblleft{}A player\textquotedblright{}.) Say it once more.
\end{tcolorbox}
\section[\texorpdfstring{\kn{ಪ್ರೇಕ್ಷಕ} (prēkṣaka)}{ಪ್ರೇಕ್ಷಕ (prēkṣaka)}]{\kn{ಪ್ರೇಕ್ಷಕ} (prēkṣaka) — a spectator, a viewer}

\label{lesson:KA-C196-preksaka}



\begin{quote}
Before the new one: say the Kannada for a director, then the Kannada for a player.
\end{quote}

\subsection*{You'll want to know: \kn{ಪ್ರೇಕ್ಷಕ}}
\textbf{\kn{ಪ್ರೇಕ್ಷಕ}} — \emph{prēkṣaka} — \textquotedblleft{}a spectator, a viewer\textquotedblright{}.

\begin{grammarlens}[title={What you've built}]
One more word for this chapter.
\end{grammarlens}

\subsection*{Your turn}
\begin{itemize}
\raggedright
  \item \emph{Say it:} \emph{prēkṣaka}
  \item \emph{Say it:} \emph{prēkṣaka}, once more
  \item \emph{Say it:} say \emph{prēkṣaka}
  \item \emph{Recall:} say the Kannada for a director, then the Kannada for a player, then say \emph{prēkṣaka} again
\end{itemize}

\begin{tcolorbox}[breakable,colback=teal!4,colframe=teal!35!black,title={Before you move on}]
What does \kn{ಪ್ರೇಕ್ಷಕ} mean? (\textquotedblleft{}A spectator, a viewer\textquotedblright{}.) Say it once more.
\end{tcolorbox}
\section[\texorpdfstring{\kn{ಬಡಗಿ} (baḍagi)}{ಬಡಗಿ (baḍagi)}]{\kn{ಬಡಗಿ} (baḍagi) — a carpenter}

\label{lesson:KA-C196-badagi}



\begin{quote}
Before the new one: say the Kannada for a player, then the Kannada for a spectator.
\end{quote}

\subsection*{You'll want to know: \kn{ಬಡಗಿ}}
\textbf{\kn{ಬಡಗಿ}} — \emph{baḍagi} — \textquotedblleft{}a carpenter\textquotedblright{}.

\begin{grammarlens}[title={What you've built}]
One more word for this chapter.
\end{grammarlens}

\subsection*{Your turn}
\begin{itemize}
\raggedright
  \item \emph{Say it:} \emph{baḍagi}
  \item \emph{Say it:} \emph{baḍagi}, once more
  \item \emph{Say it:} say \emph{baḍagi}
  \item \emph{Recall:} say the Kannada for a player, then the Kannada for a spectator, then say \emph{baḍagi} again
\end{itemize}

\begin{tcolorbox}[breakable,colback=teal!4,colframe=teal!35!black,title={Before you move on}]
What does \kn{ಬಡಗಿ} mean? (\textquotedblleft{}A carpenter\textquotedblright{}.) Say it once more.
\end{tcolorbox}
